\documentclass{article}

\usepackage{amsmath,amssymb}
\usepackage{xurl}
\usepackage[hidelinks]{hyperref}
\setlength{\emergencystretch}{2em}

% Shared commands for notes in docs/paper-gaps/.

\DeclareUnicodeCharacter{2080}{\ifmmode{}_{0}\else\textsubscript{0}\fi}
\DeclareUnicodeCharacter{2081}{\ifmmode{}_{1}\else\textsubscript{1}\fi}
\DeclareUnicodeCharacter{2082}{\ifmmode{}_{2}\else\textsubscript{2}\fi}
\DeclareUnicodeCharacter{2099}{\ifmmode{}_{n}\else\textsubscript{n}\fi}

\newcommand{\Lean}{\textsc{Lean}}
\ExplSyntaxOn
\NewDocumentCommand{\leanid}{m}{
  \ifmmode
    \textsf{\detokenize{#1}}
  \else
    \group_begin:
    \urlstyle{sf}
    \tl_set:Nn \l_tmpa_tl {#1}
    \tl_replace_all:Nnn \l_tmpa_tl {\_} {_}
    \exp_args:NV \path \l_tmpa_tl
    \group_end:
  \fi
}
\ExplSyntaxOff
\providecommand{\tr}{\operatorname{tr}}
\newcommand{\tnleanIssueBase}{https://github.com/LionSR/TNLean/issues/}
\newcommand{\tnleanPrBase}{https://github.com/LionSR/TNLean/pull/}
\newcommand{\tnleanDocBase}{https://sirui-lu.com/TNLean/docs/}
\newcommand{\ghissue}[1]{\href{\tnleanIssueBase#1}{GitHub issue \##1}}
\newcommand{\ghpr}[1]{\href{\tnleanPrBase#1}{PR \##1}}
\newcommand{\doclink}[1]{\href{\tnleanDocBase#1.html}{\path{#1}}}

% Dirac notation and the identity operator, as in blueprint/src/macros/common.tex.
% \providecommand keeps note-local definitions authoritative.
\usepackage{dsfont}
\providecommand{\ket}[1]{|#1\rangle}
\providecommand{\bra}[1]{\langle #1|}
\providecommand{\braket}[2]{\langle #1 | #2 \rangle}
\providecommand{\ketbra}[2]{|#1\rangle\!\langle #2|}
\providecommand{\Id}{\mathds{1}}
\providecommand{\one}{\mathds{1}}
\providecommand{\C}{\mathbb{C}}
\providecommand{\MN}[1]{M_{#1}(\C)}
\providecommand{\MD}{M_D(\C)}

% External referencing for paper sources.  The build helper writes sanitized
% auxiliary files under build/paper-aux/ and excludes duplicated source labels.
\usepackage{xr}

% Import a generated paper auxiliary file with a caller-supplied label prefix.
% The candidate paths support builds from docs/paper-gaps/, the repository root,
% and build/paper-aux/ itself.
\newcommand{\importpaperaux}[2]{%
  \IfFileExists{../../build/paper-aux/#2.aux}{%
    \externaldocument[#1]{../../build/paper-aux/#2}%
  }{%
    \IfFileExists{build/paper-aux/#2.aux}{%
      \externaldocument[#1]{build/paper-aux/#2}%
    }{%
      \IfFileExists{#2.aux}{%
        \externaldocument[#1]{#2}%
      }{}%
    }%
  }%
}

% General external reference macros
\newcommand{\paperref}[2]{\ref{#1-#2}}
\newcommand{\papereqref}[2]{(\ref{#1-#2})}

% CPSV16 source (1606.00608 / MPDO-22-12-17-2.tex) external document and shortcuts
\importpaperaux{cpsv16-}{1606.00608/MPDO-22-12-17-2-tnlean}
\newcommand{\cpsvref}[1]{\paperref{cpsv16}{#1}}
\newcommand{\cpsveqref}[1]{\papereqref{cpsv16}{#1}}

% Verdict marker: \gapnote{<kind>}{<status>}. Kind is the policy
% classification (clarification, local-correction, scope-restriction,
% unfaithful, false-source, open-gap); status is open, wip (elimination
% underway), resolved, or historical. Machine-read by the site index;
% prints a verdict line.
\newcommand{\gapnote}[2]{\par\noindent\textbf{Verdict:} #1 (#2).\par}


\title{Wolf's Operator-System Extension: Full-Matrix-Algebra Scope Restriction}
\author{}
\date{2026-08-06}

\begin{document}
\maketitle

\section{The source statement}

Wolf's proposition "Extending cp maps from operator systems" quantifies over
two arbitrary finite-dimensional $C^\ast$-algebras $\mathcal A,\mathcal B$: let
$T:S\to\mathcal B$ be a completely positive linear map from an operator system
$S\subseteq\mathcal A$; then there is a completely positive map
$\tilde T:\mathcal A\to\mathcal B$ agreeing with $T$ on $S$
\cite[Proposition, lines~614--626]{Wolf2012Quantum}. The local source is
\path{Notes/WolfNoteTexSource/ch01_deconstructing_quantum.tex}, lines~617--619
(statement) and 620--626 (proof). Neither $\mathcal A$ nor $\mathcal B$ is
assumed to be a full matrix algebra: a finite-dimensional $C^\ast$-algebra is,
by the finite-dimensional structure theorem, a direct sum
$\bigoplus_k M_{d_k}(\mathbb C)$, and the proposition covers every such
algebra, not only the one-summand case $\mathcal A=M_m(\mathbb C)$.

\section{The formalized scope}

The domain algebra $\mathcal A$ is now formalized at the source's stated
generality: a finite direct sum $\mathcal A=\bigoplus_{k<r} M_{d_k}(\mathbb C)$,
represented literally as the dependent function type
$\forall k, \MN{d_k}$. \leanid{Matrix.IsOperatorSystemDirectSum} and
\leanid{Matrix.IsCPOnOperatorSystemDirectSum}
(\doclink{TNLean/Channel/OperatorSystemExtensionDirectSum}) are the
direct-sum-domain operator-system and complete-positivity-at-every-level
predicates, defined entrywise (block by block), matching the single-summand
predicates \leanid{Matrix.IsOperatorSystem}/\leanid{Matrix.IsCPOnOperatorSystem}
(\doclink{TNLean/Channel/OperatorSystem}) verbatim when $r=1$. The extension
theorem \leanid{Matrix.exists_cp_extension_of_operatorSystem_directSum}
produces $\tilde T:\bigoplus_k \MN{d_k}\to\MN{p}$ agreeing with $T$ on $S$,
completely positive in the same block-by-block sense
(\leanid{Matrix.IsCPMapDirectSum}). The proof transports the single-summand
theorem \leanid{Matrix.exists_cp_extension_of_operatorSystem}
(\doclink{TNLean/Channel/OperatorSystemExtension}) along the block-diagonal
embedding $\bigoplus_k M_{d_k}(\mathbb C)\hookrightarrow M_M(\mathbb C)$,
$M=\sum_k d_k$ (\leanid{Matrix.directSumEmbedFin}, built from
\leanid{Matrix.directSumDiagonalEmbedding}, the same block-diagonal embedding
used for Wolf's Theorem~6.14 block structure, composed with the canonical
reindexing \leanid{finSigmaFinEquiv}): complete positivity at every tensor
level descends through this embedding
(\leanid{Matrix.bipartiteSlice_directSumEmbedFinLevel}), because a
block-diagonal matrix is positive semidefinite exactly when each of its
diagonal blocks is.

The codomain algebra $\mathcal B$ remains a full matrix algebra
$\mathcal B=\MN{p}$. This is not a residual scope restriction: as before, the
choice $\mathcal B=M_p(\mathbb C)$ tracks the source proof's own reduction,
which fixes the ancilla dimension and never revisits the direct-sum structure
of $\mathcal B$ after that point, so the codomain restriction is without loss
of generality for the source's own argument.

\section{Verdict}

\textbf{Closed for the domain algebra.} Wolf's proposition is now formalized
at its stated generality for the domain algebra $\mathcal A$: an arbitrary
finite direct sum $\bigoplus_{k<r} M_{d_k}(\mathbb C)$, not only the
one-summand case. No hypothesis is added beyond the source statement (the
direct-sum operator-system and complete-positivity predicates match the
source's definitions on each summand, and complete positivity at level $j$ is
exactly the source's $T\otimes\mathrm{id}_j$ positivity condition applied
block by block). The codomain restriction $\mathcal B=M_p(\mathbb C)$ remains,
documented above and in the module docstring of
\doclink{TNLean/Channel/OperatorSystemExtension} as without loss of generality
for the source's proof technique, not a gap. Tracked in \ghissue{5589}.

\bibliographystyle{alpha}
\bibliography{references}

\end{document}
